% Generated draft: requires mathematical and source review.
\documentclass[11pt]{article}
\usepackage[margin=25mm]{geometry}
\usepackage{fontspec}
\setmainfont{Latin Modern Roman}
\usepackage{amsmath,amssymb,mathtools,unicode-math}
\setmathfont{Latin Modern Math}
\usepackage{enumitem,needspace,xurl,hyperref}
\setlength{\parindent}{0pt}
\setlength{\parskip}{5pt}
\setlength{\emergencystretch}{3em}
\newcommand{\N}{\mathbb{N}}
\newcommand{\Z}{\mathbb{Z}}
\newcommand{\Q}{\mathbb{Q}}
\newcommand{\R}{\mathbb{R}}
\newcommand{\C}{\mathbb{C}}
\newcommand{\F}{\mathbb{F}}
\newcommand{\A}{\mathbb{A}}
\newcommand{\PP}{\mathbb{P}}
\newcommand{\E}{\mathbb{E}}
\newcommand{\eps}{\varepsilon}
\newcommand{\dd}{\,\mathrm{d}}
\begin{document}
\section*{UnsolvedMath ID: 700476}
\textbf{Source problem code:} catalogue entry
\subsection{Definitions and mathematical statement}
A strict total order \(<\) on a group \(G\) is invariant under multiplication on both sides if
\[
x<y\ \Longrightarrow\ axb<ayb
\qquad(a,b,x,y\in G).
\]
One standard interpretation of finite rank is finite Pr\"ufer rank: there exists an integer \(r\geq 1\) such that every finitely generated subgroup \(H\leq G\) can be generated by at most \(r\) elements. Define the derived series by
\[
G^{(0)}=G,\qquad
G^{(n+1)}=\bigl\langle u^{-1}v^{-1}uv : u,v\in G^{(n)}\bigr\rangle
\quad(n\geq 0).
\]
The group \(G\) is soluble if \(G^{(d)}=\{1\}\) for some integer \(d\geq 0\).

The supplied problem asks to prove that every linearly ordered group of finite rank is soluble. Under the interpretations specified above, its precise conditional formulation is: for every group \(G\), if \(G\) admits a strict total order invariant under multiplication on both sides and has finite Pr\"ufer rank, then
\[
\exists d\in\Z_{\geq 0}\quad G^{(d)}=\{1\}.
\]
The supplied extract does not define either the required invariance of the order or the meaning of rank. Consequently, this formulation is conditional on those conventions matching the source; their intended meanings require confirmation from the source context. No bound on \(d\) in terms of \(r\) is asserted.
\subsection{Short English statement}
V. M. Kopytov asks whether every linearly ordered group of finite rank is soluble. With a two-sided invariant order and finite Pr\"ufer rank, this means that the derived series of every such group eventually reaches the trivial subgroup. These interpretations of order and rank remain to be confirmed from the source context.
\subsection{Sources}
\begin{itemize}
\item Local UnsolvedMath registry: supplied complete source record, permanent ID \(700476\), exact source problem code \texttt{catalogue entry}, title \emph{Kourovka Notebook Problem 8.24}; upstream revision \texttt{d371ee917ac9647fbb06cd6fc529deec6f3dda82}.
\item Supplied authoritative-source attribution: \emph{The Kourovka Notebook (Unsolved Problems in Group Theory)}, original problem 8.24, 8th issue (1982), proposed by V. M. Kopytov. The record identifies the 21st issue edition (2026), \texttt{arXiv:1401.0300v48}; its collection description names E. I. Khukhro and V. D. Mazurov as editors.
\item The local record supplies the status \texttt{open}. This is a supplied registry classification, not an independently verified mathematical or literature assessment. The source edition and its terminology have not been independently checked for this draft.
\end{itemize}
\end{document}
